\documentclass{article}

\newcommand{\ta}{\mathtt{a}}
\newcommand{\tb}{\mathtt{b}}
\newcommand{\tn}{\mathtt{n}}
\newcommand{\tm}{\mathtt{m}}
\newcommand{\tx}{\mathtt{x}}
\newcommand{\ty}{\mathtt{y}}
\newcommand{\tE}{\mathtt{E}}
\newcommand{\tP}{\mathtt{P}}

\title{Interpreting attention with toy models}
\author{John MacCormick}
\date{\today}


\begin{document}

\maketitle

This document attempts to explain some aspects of the attention
mechanism in transformer neural network models. We assume the reader
is familiar with conventional definitions and treatments of the
attention mechanism. Here, we focus on some extremely simplified
variants of transformer architectures, to explain some elementary
properties of the attention mechanism that may not be obvious in fully
featured architectures.

We typically use a very small vocabulary of tokens. As an initial
example, consider the vocabulary $\{\ta, \tb, \tE, \tP\}$. In this
vocabulary and in the rest of the document: $\tE$ is the
end-of-sequence token, denoted \texttt{<EOS>} by many authors; $\tP$
is the padding token, denoted \texttt{<PAD>} by many authors; letters
such as $\ta, \tb, \tm, \tn, \tx, \ty$ denote regular tokens, which
might represent words such as \texttt{apple} or \texttt{banana} in a
more realistic application.

\end{document}
